\section{The precise conditional consequences for the padded permanent}
\label{sec:gct}

This section supplies the algebraic reductions that turn either claimed
border-circuit lower bound into a statement about determinant orbit
closures.  The reductions are unconditional; the lower bounds used in
Theorem~\ref{thm:conditional-gct} remain the claims of the two preprints.
The geometric formulation originates with Mulmuley and
Sohoni~\cite{MulmuleySohoni2001}; see also
\cite[arXiv version~3, Section~1(a), Conjecture~1.2, p.~2]
{BurgisserIkenmeyerPanova2019}.

Let $k$ be an algebraically closed field.  All closures in this section
are Zariski closures over $k$.  For $N\geq m$, put
\[
 P_{m,N}=z^{N-m}\operatorname{perm}_m,
 \qquad
 \mathcal D_N=\overline{\operatorname{GL}_{N^2}(k)\cdot\det_N}.
\]
When $N>m$, $z$ is a variable independent of the $m^2$ permanent
variables.  There is room for these $m^2+1$ variables among the $N^2$
variables on which the group acts.  When $N=m$, the factor in $z$ is
absent.  The border determinantal complexity $\underline{\operatorname{dc}}(p)$
is the least $N$ for which $p$ belongs to the coefficientwise closure of
determinants of $N$ by $N$ matrices of affine linear forms in its own
variables.  Adjoining unused variables does not change this definition:
one direction adjoins them, and the other sets them to zero.

\begin{lemma}[Independent padding]\label{lem:padding}
Let $p\in k[x_1,\ldots,x_r]$ be homogeneous of degree $m\leq N$.
If $N>m$, assume $r+1\leq N^2$ and let $z$ be independent of the
$x_i$.  Then
\[
 z^{N-m}p\in\mathcal D_N
 \quad\Longleftrightarrow\quad
 \underline{\operatorname{dc}}(p)\leq N.
\]
For $N=m$ the same assertion holds without $z$, assuming $r\leq N^2$.
The result is valid in every characteristic.
\end{lemma}

\begin{proof}
Write $\mathcal A_N$ for the set of determinants of affine linear
$N$ by $N$ matrices in the $x_i$.  Homogenization to degree $N$ is the
linear map
\[
 H_N:k[x_1,\ldots,x_r]_{\leq N}\longrightarrow
 k[x_1,\ldots,x_r,z]_N,
 \qquad H_N(q)=z^Nq(x_1/z,\ldots,x_r/z).
\]
It carries $\det(A_0+\sum_i x_iA_i)$ to
$\det(zA_0+\sum_i x_iA_i)$.  Each entry of the latter matrix is a
linear form in the ambient $N^2$ variables, so its determinant is
obtained from $\det_N$ by a possibly singular linear substitution.
Invertible endomorphisms are dense in all endomorphisms, and the map
that substitutes an endomorphism into $\det_N$ is polynomial.
Consequently every such determinant belongs to $\mathcal D_N$.
Continuity of $H_N$ now gives
$H_N(\overline{\mathcal A_N})\subseteq\mathcal D_N$.
Since $H_N(p)=z^{N-m}p$, this proves the forward implication from
border determinantal complexity to orbit closure membership.

Conversely, set $z=1$ and all unused ambient variables to zero.  This
is a linear map on coefficient spaces.  It carries every polynomial
in the orbit of $\det_N$ to a determinant of an affine linear matrix,
and carries $z^{N-m}p$ to $p$.  Continuity therefore gives
$p\in\overline{\mathcal A_N}$.  If $N=m$, replace homogenization
in the first argument by extraction of the homogeneous degree-$N$
part: that part of an affine determinant is the determinant of its
linear part.  For the converse simply set unused variables to zero.
\end{proof}

The familiar one-way padding lemma is the implication for an exact
determinantal expression.  The proof also explains why an independent
padding variable is useful: dehomogenization recovers the entire
polynomial, not a specialization of one of its variables.

\begin{lemma}[A division-free determinant circuit]\label{lem:det-circuit}
Over every commutative ring, the determinant of an $N$ by $N$ matrix
of affine linear forms in at most $N^2$ variables has a circuit with
at most $8N^4$ binary addition, subtraction and multiplication gates,
with free constants.  The same bound holds in the model with only
addition and multiplication in which every variable and constant
input is also counted as a gate.
\end{lemma}

\begin{proof}
We use the Samuelson recurrence underlying Berkowitz's algorithm;
the recurrence and its adjugate identity are
\cite[Claims~1--2, p.~148]{Berkowitz1984}.
The following elementary count suffices, without optimizing the
parallel algorithm.  For a matrix of order $n$, write
\[
 A=\begin{pmatrix}a&R\\ S&M\end{pmatrix},\qquad
 \chi_M(t)=\sum_{i=0}^{n-1}c_i t^{n-1-i},\quad c_0=1,
 \qquad b_j=RM^jS\quad(0\leq j\leq n-2).
\]
Set $c_n=0$.  Cofactor expansion of the block determinant gives
\[
 \chi_A(t)=(t-a)\chi_M(t)-R\operatorname{adj}(tI-M)S.
\]
The identity
\[
 \operatorname{adj}(tI-M)
 =\sum_{j=0}^{n-2}t^{n-2-j}
       \sum_{i=0}^{j}c_i M^{j-i}
\]
follows by multiplying by $tI-M$: all intermediate powers cancel and
the constant term is the Cayley--Hamilton identity for $M$.
These are polynomial identities over $\mathbb Z$, so they hold over
every commutative ring.  Thus, if
$\chi_A(t)=\sum_{j=0}^n d_jt^{n-j}$, its coefficients satisfy
\[
 d_0=1,\qquad
 d_j=c_j-ac_{j-1}-\sum_{i=0}^{j-2}c_i b_{j-2-i}
       \quad(1\leq j\leq n),
\]
where an empty sum is zero.

Compute the characteristic polynomial of $M$ recursively.  The
vectors $S,MS,\ldots,M^{n-2}S$ require at most
$2(n-2)(n-1)^2$ gates, and their products with $R$ at most
$2(n-1)^2$ gates.  For the coefficient recurrence, first form
$-a$ and each $-b_j$ by multiplication by the single constant $-1$.
There are $n$ such multiplications.  The formula for $d_j$, written
with these negative quantities, uses at most $j$ further
multiplications and $j$ additions.  Thus the recurrence takes
at most $n^2+2n$ gates, without a subtraction primitive.
Including the trivial order-one case and
the final sign, $4n^3$ gates per recursive order is a valid bound.
Since $\sum_{n=1}^N n^3=N^2(N+1)^2/4\leq N^4$, a generic
determinant costs at most $4N^4$ gates.  Computing each of the $N^2$
affine entries in at most $N^2$ variables costs at most $2N^2$
gates, or $2N^4$ in total.  If inputs are counted, the affine
coefficients require at most $N^2(N^2+1)$ constant inputs, the
variables at most $N^2$ inputs, and $0,1,-1$ at most three
additional inputs.  For $N\geq2$, the total is bounded by
\[
 7N^4+2N^2+3\leq8N^4.
\]
For $N=1$, a single affine form in at most one variable requires
at most five gates including inputs.  This proves both versions.
No division, interpolation or assumption on the characteristic was
used.
\end{proof}

\begin{lemma}[Homogeneous components]\label{lem:homogeneous-components}
If a circuit of size $s$ computes $q$, all homogeneous components of
$q$ of degrees at most $d$ can be computed using at most
$s(d+1)^2$ gates.  This assertion holds over every commutative ring
and also for circuits whose constants are Laurent series.
\end{lemma}

\begin{proof}
At each gate keep the list of components of degrees $0,\ldots,d$.
Addition and subtraction act componentwise.  At a multiplication
gate the component of degree $j$ is
$\sum_{i=0}^j q_i r_{j-i}$; its computation needs at most $2j+1$
gates.  The sum of these bounds for $j=0,\ldots,d$ is $(d+1)^2$.
Inputs need no gates.  Induction over the original circuit proves
the claim.  All formulas are identities over $\mathbb Z$.
\end{proof}

For completeness, the passage from geometric closure to Laurent
approximation can be made without imposing an analytic topology.

\begin{lemma}[Closure gives a Laurent approximation]
\label{lem:laurent-closure}
Let $k$ be algebraically closed and let $F:\mathbb A^a_k\to
\mathbb A^b_k$ be a polynomial map.  If
$y\in\overline{F(\mathbb A^a_k)}$, there are
$u_1(t),\ldots,u_a(t)\in k((t))$ with
$F(u(t))\in y+t k[[t]]^b$.
\end{lemma}

\begin{proof}
Close the graph of $F$ in $\mathbb P^a_k\times\mathbb A^b_k$.
Projection from this graph closure is proper, so its image is
$\overline{F(\mathbb A^a_k)}$.  Choose a closed point $x$ above
$y$.  If $x$ is on the original graph, a constant approximation
suffices.  Otherwise choose an integral curve through $x$ whose
generic point belongs to the original graph.  Here is the usual
local construction of such a curve.  On an affine neighborhood of
$x$, choose a nonzero regular function $f$ with $D(f)$ inside the
original graph.  In the local domain of dimension $d$ at $x$, extend
$f$ to a system of parameters $f,h_1,\ldots,h_{d-1}$.  A minimal
prime of $(h_1,\ldots,h_{d-1})$ has positive-dimensional quotient
by the height theorem.  It cannot contain $f$, since adjoining $f$
gives an ideal whose radical is the maximal ideal.  The quotient
has dimension one, again by the height theorem.  Its closure is
the required curve.

Normalize this curve and take a point over $x$.  Normalization is
finite for a curve of finite type over a field, and the local ring
at the chosen point is a discrete valuation ring; these are the
standard curve facts in \cite[Section~33.43, Lemmas~33.43.6 and
33.43.8]{StacksProject}.  Its residue field is $k$, since $k$ is
algebraically closed.  Choosing a uniformizer $t$ identifies its
completion with $k[[t]]$: successive subtraction of residue
coefficients gives the power-series expansion.  The source
coordinates of the graph are rational functions on the curve,
hence belong to $k((t))$.  The target coordinates are regular at
the chosen point and reduce to $y$, which proves the assertion.
\end{proof}

Apply this lemma to the map sending the coefficients of an affine
matrix to the coefficients of its determinant.  Its Laurent
approximation, followed by the division-free circuit of
Lemma~\ref{lem:det-circuit}, is an approximating circuit with the
same gate bound.  Conversely, any Laurent approximation satisfies
every polynomial equation that vanishes on the parameterized
family, by taking its constant term.  Thus this use of closure
agrees with the approximation convention of
\cite[Definitions~1.3 and N.1, pp.~2 and 62]{Anthropic2026}.
Extraction of a homogeneous component is linear on coefficient
spaces, and Lemma~\ref{lem:homogeneous-components} accounts for
its circuit cost when a bound concerns one component only.

\begin{theorem}[Conditional orbit-closure bounds]
\label{thm:conditional-gct}
Fix a characteristic $p\ne2$.  Suppose that, for every field of
characteristic $p$ and all $m\geq m_0(p)$, every approximating
circuit for $\operatorname{perm}_m$ has more than $B_p(m)$ gates.
Then, over the algebraic closure of every such field,
\begin{equation}\label{eq:gct-exact-range}
 P_{m,N}\notin\mathcal D_N
 \quad\hbox{whenever}\quad
 m\geq m_0(p),\quad N\geq m,\quad 8N^4\leq B_p(m).
\end{equation}
In particular, granting the Anthropic preprint's Theorem~1.4,
there is an absolute $a>0$ such that the conclusion holds for
\[
 m\leq N\leq
 2^{(a/8)(\log_2m)^{4/3}},
 \qquad
 m\geq m_0(p),\quad a(\log_2m)^{4/3}\geq6.
\]
Granting the OpenAI preprint's Corollary~1.2, it holds for
\[
 m\leq N\leq\exp\bigl(b_p(\log m)^{5/4}\bigr)
\]
for some $b_p>0$ and all sufficiently large $m$ of that
characteristic.  The latter conclusion applies in every
characteristic other than two, although the orbit-closure clause
of that corollary is phrased in characteristic zero.
In characteristic two, $P_{m,N}\in\mathcal D_N$ for every $N\geq m$.
\end{theorem}

\begin{proof}
Membership in $\mathcal D_N$ gives a border determinantal
expression by Lemma~\ref{lem:padding}, hence an approximating
circuit of size at most $8N^4$ by the preceding two lemmas.
This contradicts the stated lower bound whenever
\eqref{eq:gct-exact-range} holds.

For the Anthropic claim, put $L=\log_2m$ and
$B_p(m)=2^{aL^{4/3}}$, as in
\cite[Theorems~1.1 and 1.4, p.~2]{Anthropic2026}.
If $N\leq2^{(a/8)L^{4/3}}$, then
$8N^4\leq2^{3+(a/2)L^{4/3}}\leq2^{aL^{4/3}}$
when $aL^{4/3}\geq6$.  This writes out the reduction mentioned
in \cite[Remark~N.6, p.~64]{Anthropic2026}, with an explicit
loss in its constant.  The input-counting version of
Lemma~\ref{lem:det-circuit} matches the arithmetic-circuit model
in \cite[Section~1.3, pp.~3--4]{Anthropic2026}.
For the OpenAI claim, apply its
all-field border determinantal conclusion directly and decrease
the positive constant in its exponent, if necessary, to make the
displayed upper range strictly smaller than the claimed lower
bound.  Alternatively apply \eqref{eq:gct-exact-range} to its
border-circuit conclusion; a fixed fourth power only changes
that constant.  Its proof locates these reductions in
\cite[Corollary~1.2 and its proof]{OpenAI2026}.

In characteristic two the signs of permutations equal $1$,
so $\operatorname{perm}_m=\det_m$ as formal polynomials.
The block diagonal linear matrix
$\operatorname{diag}(X,zI_{N-m})$ has determinant
$z^{N-m}\det_m$.  The singular-substitution argument in
Lemma~\ref{lem:padding} places it in $\mathcal D_N$.
\end{proof}

The statements are conditional in precisely one place: the assumed
border-circuit lower bound.  Neither the reductions above nor the
paper's use of them verifies the released Lean development of
either preprint.  The Anthropic preprint explicitly excludes
Remark~N.6 from its list of formalized statements
\cite[Section~3, p.~10; Table~3, p.~67]{Anthropic2026}.

\begin{corollary}[Every fixed padding exponent]
\label{cor:fixed-padding}
Suppose the hypothesis of Theorem~\ref{thm:conditional-gct}
holds with $B_p(m)=2^{a(\log_2m)^\beta}$ for some $a>0$ and
$\beta>1$.  For every real $c\geq1$, put
$N=\lceil m^c\rceil+1$.  Then $P_{m,N}\notin\mathcal D_N$
for all sufficiently large $m$.  It suffices that $m\geq m_0(p)$
and
\[
 \log_2m\geq
 \max\left\{1,
       \left(\frac{4c+11}{a}\right)^{1/(\beta-1)}\right\}.
\]
\end{corollary}

\begin{proof}
For $m\geq2$ one has $N\leq m^c+2\leq3m^c$, and thus,
writing $L=\log_2m\geq1$,
\[
 \log_2(8N^4)\leq3+4cL+4\log_2 3
     \leq4cL+11\leq(4c+11)L.
\]
The displayed condition implies $aL^\beta\geq(4c+11)L$.
Apply Theorem~\ref{thm:conditional-gct}.
\end{proof}

For integral $c$ this includes the degree $m^c+1$ formulation of
the Mulmuley--Sohoni conjecture, and strengthens an
infinitely-many-$m$ conclusion to all sufficiently large $m$.
Both claimed bounds dominate every fixed polynomial.  They
therefore conditionally imply a cubic border determinantal lower
bound with any fixed positive constant for all sufficiently large
$m$, and also
$\underline{\operatorname{dc}}(\operatorname{perm}_m)
\geq m^2/2+1$ for all sufficiently large $m$.
An unspecified $m_0(p)$ supplies neither the stated numerical
threshold of a separate cubic theorem nor the assertion for
every $m\geq3$.  These are different quantifiers.  In particular
the finite sizes below an unknown threshold cannot be declared
checked.  The established quadratic result being compared here
is \cite[Theorem~1.0.1]{LandsbergManivelRessayre2013}.

\subsection{The stronger conditional consequence of nonflatness}

The comparison of geometric hypotheses gives a further conditional
consequence.  Its input is the OpenAI preprint's claimed stretch
theorem, rather than just its claimed unconditional separation.
The numerical constants below record the loss in passing through
a polynomial projection; they do not specify the constants in
that stretch theorem.

\begin{corollary}[Stretched-exponential conditional GCT range]
\label{cor:nonflat-stretched-gct}
Assume the price-space and spectral-duality inputs used in
Theorem~\ref{thm:crosspolytope-shortening}, and grant
\cite[Theorem~1.3]{OpenAI2026}.  If the full flat hypothesis
fails over a field $k$, then its homogeneous VNP family $f_n$
has border arithmetic circuit cost greater than $2^{n^\sigma}$
for some $\sigma>0$ and all integers $n\geq n_0$.
In particular, nonprice of the unit square suffices.

Suppose further that $\operatorname{char}k\ne2$, and choose
constants $C\geq1$ and $a>0$ such that $f_n$ is a polynomial
projection of $\operatorname{perm}_{q(n)}$, with
$q(n)\leq Cn^a$, over the algebraic closure of $k$.
Such constants exist by permanent projection completeness.
Put
\[
 \tau=\frac{\sigma}{a},\qquad
 c=(2C^{1/a})^{-\sigma},\qquad
 m_1=C\max\{2,n_0+1\}^{a}.
\]
For every integer $m\geq m_1$, the border circuit cost of
$\operatorname{perm}_m$ is greater than $2^{c m^\tau}$.
Consequently, over the algebraic closure of $k$,
\[
 P_{m,N}\notin\mathcal D_N
 \quad\hbox{if}\quad
 m\geq m_1,\quad c m^\tau\geq6,\quad
 m\leq N\leq2^{c m^\tau/8}.
\]
In particular every fixed padding
$N=\lceil m^b\rceil+1$, $b\geq1$, is excluded for all
sufficiently large $m$ under these hypotheses.
\end{corollary}

\begin{proof}
Theorem~\ref{thm:crosspolytope-shortening} turns failure of
flatness into a strict shortening in some fixed dimension.
The square case already gives one in dimension two.
Field-extension invariance in
Theorem~\ref{thm:price-characteristic} allows passage to the
algebraic closure.  The claimed stretch theorem then supplies
$\sigma$ and $n_0$.  Its arithmetic-circuit convention has free
inputs and constants and binary addition, subtraction and
multiplication, and its lower bound even concerns approximation
of the indicated homogeneous component.  We use only the
resulting lower bound for full polynomial approximation.

Permanent completeness over fields of characteristic different
from two is \cite[Theorem~2]{Valiant1979}; it supplies a
polynomial bound on the projection size.  Increase $C$ and
$n_0$ if needed to put that bound in the displayed form.
For $m\geq m_1$, set
\[
 n=\left\lfloor(m/C)^{1/a}\right\rfloor.
\]
Then $n\geq n_0$ and $q(n)\leq Cn^a\leq m$.
Moreover $(m/C)^{1/a}\geq2$, so
\[
 n\geq\tfrac12(m/C)^{1/a},\qquad
 n^\sigma\geq2^{-\sigma}C^{-\sigma/a}m^{\sigma/a}
             =c m^\tau.
\]
The block-diagonal substitution
$X\mapsto\operatorname{diag}(X,I_{m-q(n)})$ makes
$\operatorname{perm}_{q(n)}$ a projection of
$\operatorname{perm}_m$.  Compose it with the given projection
to $f_n$.  Substitution of variables by variables or field
constants does not increase the number of arithmetic operation
gates.  It is linear on coefficient spaces and hence preserves
border upper bounds.  A border circuit of cost at most
$2^{c m^\tau}$ for $\operatorname{perm}_m$ would therefore
contradict the cost strictly greater than $2^{n^\sigma}$ for
$f_n$.

Finally, if $N\leq2^{c m^\tau/8}$ and $c m^\tau\geq6$,
\[
 8N^4\leq2^{3+c m^\tau/2}\leq2^{c m^\tau}.
\]
Lemmas~\ref{lem:padding} and \ref{lem:det-circuit} give the
orbit-closure exclusion.  For a fixed $b$, the quantity
$\log_2(8(\lceil m^b\rceil+1)^4)$ is $O_b(\log m)$,
whereas $c m^\tau/\log m\longrightarrow\infty$;
this proves the last assertion.
\end{proof}

The first assertion of this corollary concerns $f_n$ in every
characteristic.  Its permanent and orbit-closure assertions
exclude characteristic two, where the explicit determinant
expression above gives containment.  All of $\sigma,n_0,C,a$
remain input-dependent.  Thus the corollary strengthens the
conditional asymptotic range without supplying an explicit
small-size threshold or establishing nonflatness itself.

\subsection{Highest weight equations and their limits}

This subsection identifies the representation-theoretic existence
statement that follows from noncontainment.  It does not turn it
into a multiplicity comparison.

\begin{proposition}[Highest weight separation]
\label{prop:highest-weight}
Let $G=\operatorname{GL}(V)$ act linearly on a finite-dimensional
complex vector space $W$, and let $X,Y\subseteq W$ be
$G$-stable closed affine cones.  If $Y\not\subseteq X$, then
there is a homogeneous highest weight polynomial $h\in\mathbb C[W]$
that vanishes on $X$ and does not vanish identically on $Y$.
If $Y=\overline{G\cdot y}$, then $h(g\cdot y)\ne0$ for some
$g\in G$.
\end{proposition}

\begin{proof}
Choose $f\in I(X)\setminus I(Y)$.  Since $X$ and $Y$ are
cones, their ideals are homogeneous, and one homogeneous
component of $f$, of some degree $d$, still has this property.
The restriction map from the finite-dimensional rational
$G$-module $I(X)_d$ to $\mathbb C[Y]_d$ is nonzero and
$G$-equivariant.  Complete reducibility over $\mathbb C$
therefore gives an irreducible submodule on which restriction
is nonzero.  Its kernel is a submodule, so restriction is
injective on that irreducible submodule.  Choose a nonzero
highest weight vector $h$ in it.  Then $h\in I(X)_d$ but
$h\notin I(Y)_d$, as required.  If $h$ vanished on the
whole orbit of $y$, it would vanish on its closure, a
contradiction.  The complete reducibility and highest weight
interpretation used here are the same representation-theoretic
framework explained in
\cite[Section~1(a)]{BurgisserIkenmeyerPanova2019}.
\end{proof}

Apply the proposition with $X=\mathcal D_N$ and
$Y=\overline{G\cdot P_{m,N}}$.  The conditional bounds give
existence of such an equation, but give no degree, weight,
coefficients or algorithm for finding it.  Nor does the
proposition require $h(P_{m,N})\ne0$ in the chosen coordinates:
its assertion is nonvanishing at some point of the orbit.

An occurrence obstruction would require a representation
type present in $\mathbb C[Y]$ and absent from $\mathbb C[X]$.
A multiplicity obstruction would require its multiplicity
in $\mathbb C[Y]$ to exceed the multiplicity in
$\mathbb C[X]$.  Neither follows from the proposition.
Indeed, let $\mathbb G_m$ act by scalar multiplication on
$\mathbb C^2$, and let $X$ and $Y$ be the two coordinate
axes.  Both are orbit closures.  The polynomial given by
the second coordinate vanishes on $X$ and not on $Y$,
and is a weight vector.  Nevertheless their graded
coordinate rings are isomorphic as $\mathbb G_m$-modules:
each degree has one copy of the same character.  Thus
even every multiplicity can agree while the cones differ.

This is consistent with the no-occurrence theorem of
B\"urgisser, Ikenmeyer and Panova: for their convention,
which pads by the existing variable $X_{11}$, every
representation type that occurs for the padded permanent
also occurs for the determinant when $N\geq m^{25}$
\cite[arXiv version~3, Theorem~1.4, p.~3]
{BurgisserIkenmeyerPanova2019}.
For our independent padding variable the same conclusion
follows, for example, when $N\geq(m+1)^{25}$.  To see this,
restrict the $(m+1)$ by $(m+1)$ matrix to
$\operatorname{diag}(z,X)$.  Its permanent becomes
$z\operatorname{perm}_m$, so its degree-$N$ padding by
the existing first entry becomes $P_{m,N}$.  A singular
linear substitution belongs to the closure of invertible
ones, and therefore the orbit closure of $P_{m,N}$ is
contained in this larger padded-permanent orbit closure.
Restriction of coordinate rings is surjective, so every
representation type of the former also occurs in the
latter.  Apply their theorem with $m+1$ in place of $m$.
Their theorem does not exclude multiplicity obstructions;
the two claimed circuit lower bounds do not construct them.
Complete reducibility was essential in the argument above,
so it has not been asserted here in positive characteristic.
